% Source: https://artofproblemsolving.com/community/c3h1985295_followup_problem_6_from_mathcounts_mini_for_january_2020
% Stable identity: aops-topic-1985295-q01
By the Pythagorean theorem $AC=\sqrt{39^2-15^2}=\sqrt{1296}=36$. Let $AD=x$, so $CD=36-x$ and $BD=45-x$.

\begin{asy}
usepackage("amsmath");
size(3.6cm);
defaultpen(fontsize(10pt));
pair A=(0,0),B=(15,0),C=(0,36),D=(0,20);
draw(A--B--C--cycle);draw(B--D);
draw((0,1.8)--(1.8,1.8)--(1.8,0));
dot(D);
label("$A$",A,SW);label("$B$",B,SE);label("$C$",C,N);label("$D$",D,W);
label("$15$",(A+B)/2,S);label("$39$",(B+C)/2,NE);label("$x$",(A+D)/2,W);label("$36-x$",(D+C)/2,W);
label("\mbox{}",(min(currentpicture,user=true)+max(currentpicture,user=true))/2,black); // keeps the text after the figure black when LaTeX typesets the labels inline
\end{asy}

Triangle $ABD$ has a right angle at $A$, so
\[
15^2+x^2=(45-x)^2\iff 225=2025-90x\iff x=20 .
\]
Hence $AD=\boxed{20}$. (Check: $BD=\sqrt{225+400}=25$ and $CD+9=16+9=25$.)
